\chapter{Related Work}\label{ch:related-work}

Building on the technical foundations introduced in the previous chapter, this chapter surveys existing research relevant to this thesis. It first examines approaches to exploring, controlling, and navigating the latent spaces of generative models, then considers visual and interactive systems designed to make such exploration accessible, and finally covers evaluation strategies for generative systems and their outputs.

\section{Latent space exploration, controllability and navigation in generative models}\label{sec:3.1}

Latent space exploration in generative models did not start with diffusion models. Early research examined whether generative models could learn representations in which distinct factors of variation were separated clearly enough to be analyzed individually. Higgins et al. investigate this question in the context of variational autoencoders. They prioritize disentangled representations, where changing one part of the learned representation affects only a single property of the output, over reconstruction quality, the accuracy with which the decoder reproduces the original input. To test this, they vary one part of the representation at a time and observe how the output changes, then analyze the results factor by factor, showing that single dimensions turn out to correspond to changes such as position, scale, or rotation. \cite{higgins2016early}
This establishes a basic exploration pattern that appears again in later work: change one part of the representation in a controlled way and inspect how the output responds.

In research on generative adversarial networks (GANs) \cite{goodfellow2014gan}, this exploration principle became more closely associated with semantic control. A central problem here is no longer whether factors can be separated, but whether meaningful attributes can be manipulated in a predictable way. InterFaceGAN addresses this problem by relating the GAN's latent representation to semantic scores, how strongly an attribute such as age or gender appears in an image measured by a separate classifier, and identifying linear subspaces within that representation that support controlled editing \cite{shen2020interfacegan}. SeFa proposes a more unsupervised approach: instead of relying on classifiers and labeled attributes, it factorizes the weights of a pretrained generator directly to find these same kinds of meaningful directions, thereby avoiding manually annotated labels and separately trained classifiers \cite{shen2021closedform}. Together, these methods illustrate a broader development from basic latent-factor traversal toward explicit and interpretable semantic manipulation.

This development is more challenging in diffusion models because, unlike many earlier generative architectures, they do not typically expose an established, semantically organized global latent interface for semantic manipulation. One approach is therefore to identify useful representation spaces within the model itself. Kwon et al. locate such a space in the U-Net bottleneck, the most compressed layer at the center of the denoising network architecture, and show that modifying its representations can produce consistent semantic changes across samples and denoising timesteps \cite{kwon2023semantic}. This indicates that control in diffusion models can be understood as navigation through internal representations rather than only as prompt modification.

Other work suggests that semantic manipulation cannot always be reduced to fixed global directions. Meaningful edits may instead depend on the local geometry of the representation space and on the current stage of the denoising process \cite{park2023riemannian}. Following this perspective, Haas et al. study interpretable directions in diffusion-model h-space, the space of the deepest feature maps of the denoiser. They find that these directions tend to capture broad, transferable changes such as pose or gender, whereas an image-specific analysis instead uncovers more localized transformations confined to a particular region, such as the eyes or mouth, which can nonetheless be transferred to edit other images in a similar way \cite{haas2024interpretable}.

Another part of research focuses less on individual semantic directions and more on the behavior of the complete path between two representations. From this perspective, the key question is not only whether two endpoints can be connected, but whether the intermediate states remain stable, coherent, and meaningful. Studies of diffusion seed space illustrate this problem clearly. Conventional interpolation methods such as Linear Interpolation and Spherical Linear Interpolation may pass through poor intermediate regions in this space, leading to low-quality images with poor semantic content at those intermediate steps \cite{samuel2023normguided}. Related instability has also been observed in diffusion latent spaces more generally, where even small additive changes can cause noticeable variations in the generated output \cite{guo2023smooth}.

This path-oriented view is also reflected in diffusion editing methods. Lee et al. gradually shift the prompt conditioning from a source prompt embedding to a target prompt embedding during denoising and use the resulting noise difference to edit only the relevant image regions, leaving the rest unchanged \cite{lee2024promptinterp}. At a larger analytical scale, Zeng et al. examine h-space representation by comparing how it shifts across large, systematically varied collections of prompts and captions, for example when a specific word is added, removed, or exchanged. Because these shifts are tied directly to the prompt text, the approach scales to uncover broader patterns, such as societal biases the model associates with particular concepts \cite{zeng2024decoding}. Taken together, these studies indicate that latent-space navigation should be evaluated not only by its endpoints, but also by the continuity and stability of the entire trajectory.

Controllability can also be achieved without directly traversing a latent space, either by modifying the conditioning signal or by introducing lightweight changes to the model parameters. One line of work treats prompt-related representations themselves as the editing space. Manipulations of text embeddings can already support controlled edits and reveal semantic directions inside the embedding representation \cite{yu2024textembedding}. This idea is extended further by adaptive conditioning methods such as DATE, which update the embedding itself throughout the denoising process rather than relying on a fixed one, as already discussed in \hyperref[ch:foundations]{Chapter~2} \cite{na2025date}. A second line of work intervenes more directly in the model parameters. Gandikota et al. introduce Concept Sliders, which represent concepts as low-rank parameter directions and enable precise, composable control over attributes in images generated by diffusion models \cite{gandikota2023conceptsliders}.

\section{Visual and interactive exploration tools for generative systems}\label{sec:3.2}

Research on visual and interactive tools for generative systems has developed largely independently from work on controllability. Instead of focusing on how a model can be steered, these systems examine how users can inspect, understand, and explore generative processes. GAN Lab is an early and influential example. It presents a browser-based interactive environment for learning and experimenting with GANs through visualizations of model structure, sample distributions, and training dynamics. Its contribution is to show that generative models become more accessible when their training process, intermediate behavior, and outputs are exposed through interactive visual representations. \cite{kahng2019ganlab}

Other works expands this idea from model understanding to the exploration of large output spaces. Dream Lens addresses the problem that generative systems can produce hundreds or thousands of output candidates that are difficult to review systematically. Its interface combines visual browsing, metadata inspection, and example-based ranking to help users navigate a large collection of generated outputs. \cite{matejka2018dreamlens}

More recent tools move closer to direct latent-space interaction. Browsing the Latent Space introduces an interactive environment for exploring latent spaces in a three-dimensional generative setting. It combines dimensionality reduction, interpolation, and visual navigation to make the latent space itself part of the user interface \cite{cabezon2023browsing}. A related example in diffusion models is the latent-space explorer developed by Zhong and Huang, which integrates custom latent operations into Stable Diffusion to support conceptual and spatial exploration within a creative workflow \cite{zhong2025latentdiffusion}. These systems transform latent space from an internal model representation into an explicit surface for interaction.

Together, these tools illustrate a broader development in generative-model interfaces. Earlier systems primarily focused on explaining model behavior and training dynamics, while later systems emphasized the organization and inspection of large output spaces. More recent approaches treat latent representations themselves as interactive environments for exploration, an idea that directly inspired the development of the NELS plugin presented in this thesis.

\section{Evaluation strategies for latent navigation and generated outputs}\label{sec:3.3}

This section reviews existing approaches for evaluating generative systems and their outputs and briefly explains what they are meant to measure. It then narrows this broader range of approaches to the metrics and evaluation procedures applied in this thesis.

Evaluation in this area addresses two related but distinct questions. The first concerns the generated images themselves, for instance whether they look plausible or match the prompt. The second concerns how a navigation method behaves across multiple steps, for instance whether a path stays smooth, drifts semantically, or becomes unstable. Since these two aspects capture different qualities of the generation process, no single universal measure covers both, and the literature therefore relies on a range of evaluation approaches.

A standard reference point in generative image evaluation is the Fréchet Inception Distance (FID), introduced by Heusel et al. to compare the distributions of generated and real image sets within a learned feature space. A lower FID score is usually interpreted as a sign of better overall realism and diversity in the generated images. FID is therefore useful when two larger image sets should be compared at a global level. However, it does not directly measure whether an interpolation trajectory is smooth, whether consecutive steps change abruptly, or whether a direction-based walk remains semantically stable. FID thus provides a useful general evaluation measure, but it does not capture the sequential behavior that is central to latent-space navigation. \cite{heusel2017ttur}

For pairwise perceptual comparison, the literature commonly relies on feature-based similarity metrics that approximate human perceptual judgment. Zhang et al. show that deep perceptual representations often correspond more closely to human judgments rather than with pixel-wise differences or conventional structural measures. LPIPS (Learned Perceptual Image Patch Similarity) compares two images using feature activations extracted from multiple intermediate layers of a pretrained neural network. The differences between these activations are combined into a single perceptual distance, so similarity is assessed in terms of visual feature representations rather than exact pixel correspondence \cite{zhang2018lpips}. This makes LPIPS particularly suitable for comparing consecutive images in a sequence and estimating the magnitude of visible change between neighboring steps. In latent-space navigation, it can therefore support the analysis of continuity, abrupt transitions, and local instability.

Another class of metrics evaluates semantic alignment between text and image. CLIPScore estimates how well an image corresponds to a textual description by measuring the cross-modal similarity between their embeddings, without requiring a reference image \cite{hessel2022clipscore}. ImageReward extends this approach by training a reward model on human preference data for text-to-image generation. It combines text and image features from BLIP (Bootstrapping Language-Image Pre-training), a pretrained vision-language model, through cross-attention and maps the resulting representation to a scalar reward score using a multilayer perceptron, with higher values indicating outputs that are more likely to align with human preferences \cite{xu2023imagereward}. These metrics are particularly useful when evaluating whether a generated image remains consistent with an intended prompt or semantic concept.

At the same time, the literature also shows that automatic metrics alone are not enough. Otani et al. show that commonly used measures such as FID and CLIPScore may not fully align with human perception and can become increasingly difficult to interpret as model quality improves. They therefore propose a standardized human-evaluation protocol in which individual text-image pairs are assessed according to criteria such as visual fidelity and text-image alignment \cite{otani2023human}. Their work highlights the importance of structured qualitative assessment as a complement to automatic metrics, particularly when continuity, artifacts, semantic plausibility, and interpretability must be considered together.

The present thesis does not apply every metric discussed above, as the available evaluation methods cover a broader scope than this work can reasonably address. Its objective is not to benchmark text-to-image models in general, but to evaluate a latent-space exploration plugin and the results produced by its workflow. The evaluation therefore concentrates on methods that correspond directly to the system's outputs and research questions. These include neighborhood- and cluster-based measures, such as k-nearest-neighbor agreement, within-cluster similarity, and the proportion of points classified as noise. Trajectory behavior is assessed using LPIPS together with structured qualitative inspection of generated images and their corresponding graph trajectories. This is a methodological choice for the focused scope of this thesis, because these measures make it possible to assess the plugin's graph structure and navigation outputs directly, whereas broad model-level benchmark scores address a different evaluation objective.
